% Licensed under the Solderpad Hardware License v 2.1
% See: https://solderpad.org/licenses/SHL-2.1/

\chapter{Performance Optimized Architecture: Domain-Specific Architectures and Accelerators}\label{chap:dsa_architecture}
\index{domain-specific architecture (DSA)}\index{accelerator: specialized hardware compute unit}\index{performance optimization!architecture level}
\localtoc

\section{The Post-ILP Era and the Rise of Domain-Specific Architectures}\label{sec:post_ilp_dsa}
\index{Dennard scaling}\index{power wall}\index{Amdahl's law}\index{Hennessy, John L.}\index{Patterson, David A.}

The preceding chapters established two foundational tiers of processor performance optimization:
\begin{enumerate}
    \item \textbf{Digital Logic Optimization (Chapter \ref{chap:perf_opt_logic})}: Enhancing the underlying circuit building blocks (prefix adders, multi-ported register files, high-speed multipliers, and critical-path timing closure) to minimize gate delay and maximize clock frequency ($f_{\text{clock}} = 1/T_{\text{clock}}$).
    \item \textbf{Organization and Microarchitecture Optimization (Chapter \ref{chap:perf_opt_org})}: Reorganizing instruction execution through pipelining, branch prediction, superscalar dispatch, dynamic out-of-order scheduling, and hardware multithreading to minimize Clock Cycles per Instruction ($\text{CPI}$).
\end{enumerate}

For several decades, these circuit and microarchitectural techniques fueled exponential growth in general-purpose computing performance, driven by \textbf{Moore's Law} (doubling transistor density every 18--24 months) and \textbf{Dennard scaling} (constant power density through proportional voltage and capacitance scaling). Under Dennard scaling, each new semiconductor process generation enabled smaller, faster transistors that operated at higher clock frequencies without increasing power consumption per unit area.

\subsection{The Triad of Physical Limits}
Around 2005, this classical scaling paradigm collided with three fundamental physical limits:
\begin{enumerate}
    \item \textbf{The Breakdown of Dennard Scaling and the Power Wall}: As transistor channel lengths shrank below $90\,\text{nm}$, sub-threshold leakage current escalated dramatically. Supply voltages ($V_{dd}$) could no longer be scaled downward without causing thermal instability and bit errors. Consequently, power density soared, creating the \emph{Power Wall}. Processors could no longer sustain escalating clock frequencies, forcing operating frequencies to plateau in the $3\text{--}5\,\text{GHz}$ range and leaving dark silicon (portions of the die that must remain unpowered to avoid thermal destruction).
    \item \textbf{The ILP Wall}: Expanding instruction issue width (e.g., from 4-way to 8-way or 16-way superscalar) yielded severely diminishing performance returns. Real-world serial programs contain finite instruction-level parallelism bounded by true data dependencies, irregular control flow, and memory pointer aliasing.
    \item \textbf{The Memory Wall}: Microprocessor performance historically improved at $\sim 55\%$ per year, while DRAM memory latency improved at only $\sim 7\%$ per year. This divergence created a multi-hundred-cycle penalty for off-chip memory access, causing deeply pipelined superscalar cores to spend most of their execution time stalled awaiting memory operands.
\end{enumerate}

\subsection{The Overhead of General-Purpose Out-of-Order Cores}
In a modern out-of-order superscalar processor, only a tiny fraction of total die area and energy is devoted to actual arithmetic computation (the ALUs). Over $80\text{--}90\%$ of the energy and silicon footprint is consumed by \textbf{von Neumann control overhead}:
\begin{itemize}
    \item Fetching and decoding 32-bit or 64-bit instruction streams.
    \item Managing multi-level branch target buffers and pattern history tables.
    \item Dynamic register renaming tables and physical register mapping.
    \item Associative reservation station wakeup and select CAM logic.
    \item Reorder buffer FIFO retirement and precise exception tracking.
\end{itemize}
When executing an 8-bit or 16-bit arithmetic operation, a general-purpose CPU may expend several nanojoules of energy in control overhead just to deliver a few picojoules of actual arithmetic computation.

\subsection{Domain-Specific Architectures (DSAs) and Amdahl's Law}
As recognized by John L. Hennessy and David A. Patterson in their ACM Turing Lecture \citep{hennessy2019new}, the primary path forward for computer architecture is the development of \textbf{Domain-Specific Architectures (DSAs)}. Rather than attempting to accelerate all arbitrary software workloads equally, a DSA tailors its instruction set architecture, register storage, and memory dataflow to a specific class of computational problems (such as scientific computing, vector processing, digital signal processing, or neural network inference).

By specializing the hardware architecture to the domain:
\begin{itemize}
    \item \textbf{Control Overhead is Amortized}: A single instruction triggers hundreds or thousands of arithmetic operations (as in vector processors and matrix engines), reducing instruction fetch and decode energy to negligible levels.
    \item \textbf{Memory Dataflow is Optimized}: Domain-specific memory hierarchies (streaming buffers, scratchpad memories, and systolic data reuse) minimize expensive DRAM accesses.
    \item \textbf{Arithmetic Precision is Tailored}: Hardware datapaths eliminate unnecessary floating-point mantissa bits, utilizing narrow integers (INT8, INT4) or specialized micro-floating-point formats (BF16, FP8) to pack tens of times more multipliers into the same silicon area.
\end{itemize}

According to \textbf{Amdahl's Law}, the overall speedup achieved by an accelerator is strictly limited by the fraction of the program that cannot be accelerated ($s$):
\begin{equation}
\text{Speedup}_{\text{overall}} = \frac{1}{(1 - f) + \frac{f}{\text{Speedup}_{\text{accelerator}}}}
\end{equation}
where $f$ is the fraction of total execution time spent in the specialized computational kernel. Consequently, practical modern systems are organized as \textbf{heterogeneous architectures}, coupling an efficient general-purpose host CPU (managing operating system tasks, I/O, and irregular serial code) with one or more specialized domain accelerators.

\subsection{Framing DSAs via the RISC-V Modular Architecture}
The open-standard \textbf{RISC-V} architecture \citep{asanovic2014instruction, riscv-spec} provides the ideal academic and industrial vehicle for exploring domain-specific architectures. Rather than enforcing an immutable monolithic instruction set, RISC-V is architected from first principles as:
\begin{enumerate}
    \item A frozen, ultra-minimalist base integer instruction set (\texttt{RV32I} or \texttt{RV64I}) consisting of just 40 instructions.
    \item A library of standard, orthogonal extensions tailored to distinct computational domains:
    \begin{itemize}
        \item \textbf{Standard Extension `F' and `D'}: IEEE 754 Single- and Double-Precision Floating-Point.
        \item \textbf{Standard Extension `V'}: Scalable, length-agnostic Vector Processing.
        \item \textbf{Standard Extension `P'}: Packed-SIMD for Digital Signal Processing (DSP).
        \item \textbf{Standard Extension `Zfh'}: Half-precision floating-point (FP16 and BF16) for AI.
        \item \textbf{Matrix Extensions (RV-M / Zve)}: 2D tile matrix multiplication for deep learning.
    \end{itemize}
\end{enumerate}
In the following sections, we explore how each computational domain transforms processor architecture to maximize computational density, bandwidth, and energy efficiency.


\section{High-Performance Computing and Floating-Point Architecture}\label{sec:hpc_floating_point}
\index{floating-point unit (FPU)}\index{IEEE 754 standard}\index{fused multiply-add (FMA)}\index{scientific computing}

High-Performance Computing (HPC), physical simulations, weather forecasting, computational fluid dynamics, and molecular dynamics require arithmetic operations on real numbers across enormous dynamic ranges. This computational domain is fundamentally powered by \textbf{Floating-Point Units (FPUs)} conforming to the IEEE 754 standard (detailed in Chapter \ref{chap:floating_point}).

\subsection{Architectural Requirements of IEEE 754 Arithmetic}
Unlike fixed-point integer arithmetic, IEEE 754 floating-point numbers represent numbers as $(-1)^s \times (1 + \text{fraction}) \times 2^{\text{exponent} - \text{bias}}$. Supporting this standard imposes significant microarchitectural requirements:
\begin{itemize}
    \item \textbf{Alignment and Normalization Shifters}: Floating-point addition requires aligning exponents by shifting the smaller mantissa, followed by post-addition normalization and leading-zero detection.
    \item \textbf{Rounding Logic}: The architecture must support five standardized rounding modes: Round to Nearest Ties to Even (\texttt{RNE}), Round towards Zero (\texttt{RTZ}), Round Down towards $-\infty$ (\texttt{RDN}), Round Up towards $+\infty$ (\texttt{RUP}), and Round to Nearest Ties to Max Magnitude (\texttt{RMM}).
    \item \textbf{Exception Flags}: The architecture must record accrued exception bits: Inexact (\texttt{NX}), Underflow (\texttt{UF}), Overflow (\texttt{OF}), Divide by Zero (\texttt{DZ}), and Invalid Operation (\texttt{NV}).
\end{itemize}

\subsection{The RISC-V Floating-Point Architecture (`F' and `D')}
In the RISC-V ISA \citep{riscv-spec}, floating-point capability is added through the standard extensions:
\begin{itemize}
    \item \textbf{RV32F / RV64F}: Single-precision (32-bit) IEEE 754 floating-point.
    \item \textbf{RV32D / RV64D}: Double-precision (64-bit) IEEE 754 floating-point.
\end{itemize}

Architecturally, RISC-V decouples floating-point operations from the integer core through several key design decisions:
\begin{enumerate}
    \item \textbf{Dedicated Floating-Point Register File}: The architecture provides 32 independent floating-point registers, \texttt{f0} through \texttt{f31}. In an \texttt{RV32D} core, each floating-point register is 64 bits wide. Unlike the integer register file where \texttt{x0} is hardwired to zero, \textbf{register \texttt{f0} is a full, readable/writable register}. Single-precision values stored in 64-bit registers use \emph{NaN-boxing}, filling the upper 32 bits with 1s to detect typing errors.
    \item \textbf{Floating-Point Control and Status Register (\texttt{fcsr})}: A dedicated 32-bit CSR containing:
    \begin{itemize}
        \item Bits \texttt{[4:0]} (\texttt{fflags}): Accrued exception flags (\texttt{NV}, \texttt{DZ}, \texttt{OF}, \texttt{UF}, \texttt{NX}).
        \item Bits \texttt{[7:5]} (\texttt{frm}): Dynamic rounding mode selector.
    \end{itemize}
    \item \textbf{Fused Multiply-Add (FMA)}: The architecture includes four-operand FMA instructions (\texttt{fmadd.s}, \texttt{fmsub.s}, \texttt{fnmadd.s}, \texttt{fnmsub.s}) that compute:
    \begin{equation}
    d = (a \times b) \pm c
    \end{equation}
    with \textbf{only a single rounding error} at the final step. FMA doubles the floating-point execution throughput (performing two floating-point operations, FLOPs, per instruction) while dramatically improving numerical precision in dot products and matrix decompositions.
\end{enumerate}

\subsection{Pipelined Floating-Point Execution Datapaths}
Because floating-point operations involve complex multi-step algorithms (exponent subtraction, alignment barrel shifting, mantissa multiplication, normalization, leading-zero anticipation, and rounding), floating-point functional units require multi-cycle pipelining:
\begin{itemize}
    \item \textbf{Pipelined FP Adder}: Typically structured in 3 to 4 pipeline stages (Align, Add, Normalize, Round).
    \item \textbf{Pipelined FP Multiplier / FMA}: Typically structured in 4 to 5 pipeline stages (Booth recoding, Wallace/Dadda tree partial product compression, carry-propagate addition, normalization, and rounding).
    \item \textbf{Iterative FP Divider / Square Root}: Because non-restoring or Goldschmidt division is iterative, division is typically unpipelined or multi-cycle, executing over 10 to 30 clock cycles while non-blocking reservation stations allow independent integer and floating-point instructions to proceed.
\end{itemize}


\section{Vector Processing and Data-Level Parallelism}\label{sec:vector_processing}
\index{vector processor}\index{data-level parallelism (DLP)}\index{RISC-V!vector extension (RVV)}\index{SIMD vs. vector}\index{Cray, Seymour}

Many scientific, machine learning, multimedia, and graphics algorithms exhibit massive \textbf{Data-Level Parallelism (DLP)}: identical mathematical operations performed across large arrays, vectors, or streams of independent data elements.

\subsection{Evolution: Fixed-Width Packed SIMD vs. Scalable Vectors}
Historically, computer architects developed two distinct approaches to exploit data-level parallelism:

\subsubsection{1. Fixed-Width Packed SIMD (MMX, SSE, AVX, NEON)}
In the late 1990s, general-purpose microprocessor vendors introduced packed Single Instruction, Multiple Data (SIMD) extensions \citep{conte1996simd, firasta2008intel}:
\begin{itemize}
    \item Standard integer registers (or newly added SIMD registers) were partitioned into fixed-width sub-words. For example, a 128-bit register was partitioned into four 32-bit integers, with a single opcode instructing four parallel ALUs to perform four simultaneous additions.
    \item \textbf{Architectural Flaws of Packed SIMD}:
    \begin{itemize}
        \item \emph{Opcode Explosion}: Every time the hardware register width doubled (from 64-bit MMX to 128-bit SSE, 256-bit AVX, and 512-bit AVX-512), the architecture introduced hundreds of new opcodes.
        \item \emph{Binary Incompatibility}: Software compiled for 256-bit AVX2 cannot run on a 128-bit core, and cannot automatically utilize the 512-bit execution units of an AVX-512 processor without complete recompilation and algorithmic retuning.
        \item \emph{Loop-Ending Cleanup}: If a loop processes an array of length $N = 103$ elements, a 4-wide SIMD processor must execute 25 vector iterations followed by a sequential scalar ``cleanup loop'' to process the remaining 3 elements.
    \end{itemize}
\end{itemize}

\subsubsection{2. True Scalable Vector Architectures (Cray-1 to RVV)}
Pioneered by Seymour Cray in the **CRAY-1** supercomputer in 1976 \citep{cray1978cray1}, true vector architectures operate on deep vectors of configurable length using vector registers and pipelined execution lanes:
\begin{itemize}
    \item The instruction specifies operations on entire one-dimensional mathematical vectors, leaving the hardware free to execute the vector across whatever physical width (lanes) the silicon budget permits.
    \item Modern vector architecture is embodied in the **RISC-V Vector Extension (RVV)** \citep{riscv-spec}, which adopts a **Vector-Length Agnostic (VLA)** design.
\end{itemize}

\subsection{The RISC-V Vector Extension (RVV) Architecture}
The RVV architecture provides 32 vector registers, \texttt{v0} through \texttt{v31}. Its operational semantics are governed by dynamic hardware control registers:
\begin{itemize}
    \item \textbf{\texttt{VLEN}}: The hardware implementation parameter defining the exact length of each physical vector register in bits (e.g., $\text{VLEN} = 128$ for embedded IoT, $\text{VLEN} = 512$ for desktop/server, $\text{VLEN} = 2048$ for supercomputing).
    \item \textbf{\texttt{SEW} (Selected Element Width)}: The dynamic bit width of individual elements (e.g., 8, 16, 32, or 64 bits).
    \item \textbf{\texttt{LMUL} (Vector Register Multiplier)}: Allows multiple adjacent vector registers to be grouped into a single logical vector ($LMUL = 1, 2, 4, 8$) to increase vector capacity and amortize loop overhead.
    \item \textbf{\texttt{vl} (Vector Length)}: An architectural register that specifies the exact number of active elements to be processed by subsequent vector instructions.
\end{itemize}

\paragraph{Dynamic Vector Configuration: \texttt{vsetvli}}
The heart of the RVV programming model is the \texttt{vsetvli} (Vector Set Vector Length Immediate) instruction:
\begin{verbatim}
    vsetvli rd, rs1, vtype
\end{verbatim}
The programmer provides the total number of remaining array elements in register \texttt{rs1}. The hardware dynamically evaluates \texttt{rs1} against its internal $\text{VLMAX} = (\text{VLEN} / \text{SEW}) \times \text{LMUL}$, sets the architectural vector length register \texttt{vl} to $\min(\text{rs1}, \text{VLMAX})$, and returns this value in \texttt{rd}.

\paragraph{Stripmining Without Cleanup Loops}
This dynamic mechanism enables elegant \textbf{stripmining} in software:
\begin{verbatim}
# Vector Addition: C[i] = A[i] + B[i] for N elements (32-bit floats)
# Inputs: a0 = N, a1 = ptr A, a2 = ptr B, a3 = ptr C
loop:
    vsetvli t0, a0, e32, m1, ta, ma   # Configure for 32-bit floats; t0 = vl
    vle32.v  v1, (a1)                  # Load vl elements from array A
    vle32.v  v2, (a2)                  # Load vl elements from array B
    vfadd.vv v3, v1, v2                # Vector-vector floating-point add
    vse32.v  v3, (a3)                  # Store vl elements into array C
    slli    t1, t0, 2                 # t1 = t0 * 4 bytes per element
    add     a1, a1, t1                # Advance pointer A
    add     a2, a2, t1                # Advance pointer B
    add     a3, a3, t1                # Advance pointer C
    sub     a0, a0, t0                # Decrement remaining count N by vl
    bnez    a0, loop                  # Repeat until all elements processed
\end{verbatim}
This exact binary code runs unmodified on any hardware implementation—whether a low-power core with a 128-bit vector datapath or a high-performance datacenter processor with 1024-bit vector registers—always operating at 100\% hardware efficiency without a single scalar cleanup instruction!

\subsection{Vector Execution Lanes and Memory Operations}
To execute vector instructions efficiently, a vector core partitions its datapath into parallel **lanes**, as shown in Figure \ref{fig:vector_lanes}:
\begin{figure}[htbp]
  \begin{center}
    \small
    \begin{tabular}{|c|c|c|c|}
      \hline
      \textbf{Lane 0} & \textbf{Lane 1} & \textbf{Lane 2} & \textbf{Lane 3} \\ \hline
      Vector Reg Slice 0 & Vector Reg Slice 1 & Vector Reg Slice 2 & Vector Reg Slice 3 \\ \hline
      Integer / FP ALU 0 & Integer / FP ALU 1 & Integer / FP ALU 2 & Integer / FP ALU 3 \\ \hline
      Memory Port 0 & Memory Port 1 & Memory Port 2 & Memory Port 3 \\ \hline
    \end{tabular}
    \caption{Microarchitecture of a 4-lane vector processor. Each lane contains a proportional slice of the vector register file, execution ALUs, and memory interconnect.}
    \label{fig:vector_lanes}
  \end{center}
\end{figure}

RVV supports three fundamental classes of vector memory access:
\begin{enumerate}
    \item \textbf{Unit-Stride}: Sequential memory access ($A[i], A[i+1], A[i+2]$), maximizing cache line burst bandwidth.
    \item \textbf{Strided}: Accessing elements separated by a fixed stride step ($A[i \times s]$), vital for multi-dimensional array columns.
    \item \textbf{Indexed (Gather / Scatter)}: Accessing elements at irregular addresses specified by an index vector ($A[\text{index}[i]]$), essential for sparse matrix operations and graph analytics.
\end{enumerate}


\section{Graphics Processing Units (GPUs) and Massively Parallel SIMT}\label{sec:gpu_simt}
\index{graphics processing unit (GPU)}\index{single instruction, multiple threads (SIMT)}\index{warp}\index{branch divergence}

While vector architectures exploit data-level parallelism through a single thread issuing wide vector instructions, **Graphics Processing Units (GPUs)** take a distinct architectural path: **Single Instruction, Multiple Threads (SIMT)** \citep{owens2008gpu}.

\subsection{SIMD vs. SIMT Programming Models}
\begin{itemize}
    \item In **SIMD/Vector**, the programmer writes code that explicitly manipulates vector registers containing multiple data elements.
    \item In **SIMT**, the programmer writes scalar code designed to execute on a \emph{single thread}. The hardware then spawns thousands of concurrent threads and groups them into lockstep bundles called **warps** (in NVIDIA terminology) or **wavefronts** (in AMD terminology), typically 32 threads wide.
\end{itemize}
A single instruction fetch and decode unit broadcasts a common instruction across all 32 execution units in the warp, achieving SIMD-like execution density while presenting a scalar thread programming interface to software.

\subsection{Hiding Latency Through Massive Thread Scheduling}
Unlike CPUs, which allocate massive silicon area to deep out-of-order execution windows, multi-level branch predictors, and large caches to minimize instruction latency, GPUs are fundamentally **latency-hiding, throughput-oriented engines**:
\begin{itemize}
    \item A GPU core maintains dozens of warps (thousands of active threads) in hardware simultaneously, with their full architectural state held in an immense register file (often $256\,\text{KB}$ per Streaming Multiprocessor).
    \item When Warp 0 initiates an off-chip global memory read that will take 400 clock cycles to complete, the warp scheduler does not stall the pipeline. It immediately switches in \textbf{zero clock cycles} to Warp 1, Warp 2, or Warp 3.
    \item By interleaving hundreds of threads across pipelined execution units, memory latency is completely masked by arithmetic computation.
\end{itemize}

\subsection{The Challenge of Branch Divergence}
The major vulnerability of the SIMT model is **branch divergence**:
\begin{verbatim}
    if (threadIdx.x < 16) {
        // Path A
        result = compute_A();
    } else {
        // Path B
        result = compute_B();
    }
\end{verbatim}
When different threads within the same 32-thread warp evaluate a conditional branch differently:
\begin{enumerate}
    \item The warp execution unit must \textbf{serialize} the execution paths.
    \item First, it executes Path A while setting an internal hardware execution mask to disable threads 16 through 31.
    \item Next, it executes Path B while inverting the execution mask to disable threads 0 through 15.
    \item Finally, execution reconverges at the join point.
\end{enumerate}
During branch divergence, execution throughput drops by 50\% or more, illustrating why GPUs achieve maximum efficiency on uniform, regular compute kernels.

\subsection{Open-Source RISC-V GPGPU Architectures}
The principles of SIMT execution have recently been applied to open-source hardware through architectures such as **Vortex**, a RISC-V GPGPU. Vortex augments the base RISC-V ISA with hardware warp scheduling, SIMT stack convergence instructions, and high-bandwidth shared scratchpad memories, enabling native execution of OpenCL and Vulkan compute workloads on an open-standard RISC-V foundation.


\section{Digital Signal Processors (DSPs) and Embedded Compute}\label{sec:dsp_architecture}
\index{digital signal processor (DSP)}\index{fixed-point arithmetic}\index{multiply-accumulate (MAC)}\index{saturating arithmetic}\index{RISC-V!packed SIMD (P)}

Digital Signal Processors (DSPs) represent one of the earliest and most successful classes of domain-specific architectures \citep{eyre2000evolution}. Tailored for real-time streaming data—such as telecommunications (5G/LTE physical layers), audio processing, radar, and sensor telemetry—DSPs optimize for continuous, deterministic mathematical transformations.

\subsection{Characteristics of Signal Processing Workloads}
Signal processing kernels are dominated by two mathematical operations:
\begin{enumerate}
    \item **Finite Impulse Response (FIR) Convolution and Correlation**:
    \begin{equation}
    y[n] = \sum_{k=0}^{M-1} h[k] \cdot x[n-k]
    \end{equation}
    \item **Fast Fourier Transforms (FFT)**: Decomposing time-domain signals into frequency spectra via recursive butterfly computations:
    \begin{equation}
    X[k] = A + W_N^k \cdot B
    \end{equation}
\end{enumerate}
Both kernels require millions of back-to-back \textbf{Multiply-Accumulate (MAC)} operations executed on continuous data streams.

\subsection{Fixed-Point Arithmetic and the $Q$-Format}
While scientific simulations rely on floating-point arithmetic, embedded real-time signal processing historically avoided floating-point due to silicon area, latency, and power constraints. Instead, DSPs utilize **fixed-point arithmetic**, typically formatted as $Q$-numbers:
\begin{itemize}
    \item A signed $Q1.15$ number represents a fractional value in the range $[-1.0, +1.0 - 2^{-15}]$ using a 16-bit two's complement integer, where the most significant bit is the sign bit and the remaining 15 bits represent fractional binary digits ($2^{-1}, 2^{-2}, \dots, 2^{-15}$).
    \item When two $Q1.15$ numbers are multiplied, the product is a 32-bit $Q2.30$ number, which must be shifted left by 1 bit to restore the $Q1.31$ format.
\end{itemize}

\subsection{Hardware Architectural Features of DSPs}
To execute real-time signal processing with maximum efficiency, DSP architectures incorporate five distinct features:
\begin{enumerate}
    \item \textbf{Single-Cycle Multiply-Accumulate (MAC) Units}: A hardware datapath containing a fast multiplier cascaded into an adder with an accumulator register ($A = A + (X \times Y)$), equipped with \emph{guard bits} (e.g., a 40-bit accumulator for 32-bit products) to prevent intermediate overflow during iterative summations.
    \item \textbf{Saturating Arithmetic}: In standard integer arithmetic, overflow causes two's complement wrap-around (e.g., $+32767 + 1 = -32768$), which introduces catastrophic audible popping or control instability in signal processing. In contrast, DSP arithmetic units provide \textbf{saturation}: when an arithmetic result exceeds the maximum representable bound, the value is clamped to the maximum positive ($0\text{x}7\text{FFF}$) or negative ($0\text{x}8000$) value.
    \item \textbf{Modified Harvard Architecture with Dual Memory Access}: Digital signal processing algorithms require fetching an instruction, an input data sample $x[n-k]$, and a filter coefficient $h[k]$ simultaneously in a single clock cycle. Classical DSPs (such as the Texas Instruments TMS320 series) provide two or three distinct internal memory banks and buses ($X$-data bus, $Y$-data bus, and Program bus).
    \item \textbf{Circular Buffer Addressing}: FIR filters model sliding data windows. Hardware Address Generation Units (AGUs) automatically wrap memory pointers around circular buffers of size $M$ without software branch instructions or modulo operations.
    \item \textbf{Bit-Reversed Addressing}: In-place Decimation-in-Time FFT algorithms produce output samples in bit-reversed order (e.g., index $001_2 \leftrightarrow 100_2$). Specialized DSP AGUs automatically reverse index bits during address calculation, eliminating memory permutation passes.
\end{enumerate}

\subsection{The RISC-V Packed-SIMD / DSP Extension (`P')}
In modern RISC-V systems, DSP capability is integrated via the standard **`P' Extension (Packed-SIMD)**. Rather than introducing distinct DSP registers, the `P` extension partitions existing 32-bit or 64-bit general-purpose integer registers into parallel sub-word vectors (e.g., dual 16-bit or quad 8-bit integers). It provides single-cycle saturating addition, fractional multiplication, and dual-MAC instructions, delivering high-performance audio and communications processing on low-power embedded cores without the silicon area cost of a dedicated vector unit.


\section{Tensor and Matrix Processors for Artificial Intelligence}\label{sec:tensor_matrix_processors}
\index{tensor processing unit (TPU)}\index{systolic array}\index{artificial intelligence}\index{matrix multiplication}\index{Jouppi, Norman P.}

In the 2010s, the explosive emergence of deep learning transformed computer architecture. Deep neural networks (convolutional networks, Transformers, and Large Language Models) exhibit compute requirements that dwarf traditional software workloads, growing at over $5\times$ to $10\times$ per year.

\subsection{Mathematical Core: General Matrix Multiply (GEMM)}
Over 90\% of the computational effort in modern neural network training and inference consists of **General Matrix Multiply (GEMM)** operations:
\begin{equation}
C = \alpha (A \times B) + \beta C
\end{equation}
Multiplying two $N \times N$ matrices requires $O(N^3)$ arithmetic multiply-accumulate operations operating on only $O(N^2)$ data elements. This high ratio of arithmetic operations to memory bytes—known as **Arithmetic Intensity** (FLOPs per byte transferred)—means that matrix multiplication is intrinsically compute-bound rather than memory-bound, provided that data can be reused effectively on-chip.

\subsection{Systolic Arrays: Amortizing Memory Access Through Spatial Reuse}
In traditional CPU and GPU architectures, every multiplication requires fetching operands from a central register file and writing the result back through cache hierarchies. For high-dimensional matrix multiplication, this register-file communication consumes substantial energy.

In 1982, H. T. Kung introduced the **Systolic Array** architecture. Like the biological circulatory system pumping blood through tissue, a systolic array pumps data operands rhythmically across a 2D grid of tightly coupled Processing Elements (PEs), as depicted in Figure \ref{fig:systolic_array}.

\begin{figure}[htbp]
  \begin{center}
    \small
    \begin{tabular}{|c|c|c|c|}
      \hline
      \textbf{Input $A$} $\to$ & \textbf{PE}$_{0,0}$ $\to$ & \textbf{PE}$_{0,1}$ $\to$ & \textbf{PE}$_{0,2}$ $\to$ \\ \hline
      \textbf{Input $A$} $\to$ & \textbf{PE}$_{1,0}$ $\to$ & \textbf{PE}$_{1,1}$ $\to$ & \textbf{PE}$_{1,2}$ $\to$ \\ \hline
      \textbf{Input $A$} $\to$ & \textbf{PE}$_{2,0}$ $\to$ & \textbf{PE}$_{2,1}$ $\to$ & \textbf{PE}$_{2,2}$ $\to$ \\ \hline
      & $\downarrow$ \textbf{Data $B$} & $\downarrow$ \textbf{Data $B$} & $\downarrow$ \textbf{Data $B$} \\ \hline
    \end{tabular}
    \caption{Conceptual 2D systolic array. Matrix operands stream rhythmically across horizontal and vertical interconnects between neighboring Processing Elements (PEs).}
    \label{fig:systolic_array}
  \end{center}
\end{figure}

In a classic **Weight-Stationary** systolic array (such as the Google Tensor Processing Unit, TPU v1 \citep{jouppi2017indatacenter}):
\begin{enumerate}
    \item Weights from matrix $B$ are loaded once and held stationary inside the local registers of each PE.
    \item Activation values from matrix $A$ stream horizontally across the array from left to right.
    \item Partial sum accumulations stream vertically from top to bottom.
    \item Each PE computes a local multiply-accumulate operation ($P_{\text{out}} = P_{\text{in}} + A \times B$) in a single cycle and passes the activation directly to its right neighbor without touching the memory bus.
\end{enumerate}
A $256 \times 256$ systolic array executes $65,536$ multiply-accumulate operations (131,072 arithmetic operations) every single clock cycle, while accessing off-chip memory only at the array perimeter!

\subsection{Low-Precision AI Numerics}
Deep neural networks exhibit remarkable resilience to numerical noise. High-precision 64-bit or 32-bit floating-point numbers are completely unnecessary for neural network inference.

Architectures exploit this by adopting narrow, energy-efficient numerical formats:
\begin{itemize}
    \item \textbf{INT8 and INT4}: Quantizing weights and activations to 8-bit or 4-bit integers reduces memory bandwidth by $4\times$ to $8\times$, while an 8-bit multiplier occupies $\sim 1/16^{\text{th}}$ the silicon area and consumes $\sim 1/20^{\text{th}}$ the energy of a 32-bit floating-point multiplier.
    \item \textbf{Bfloat16 (Brain Floating Point)}: Developed by Google, Bfloat16 preserves the full 8-bit dynamic range of IEEE 754 single-precision float while reducing the mantissa to 7 bits, enabling direct drop-in acceleration during neural network training without complex underflow scaling.
    \item \textbf{FP8 Formats (E4M3 and E5M2)}: Emerging 8-bit floating-point standards providing high dynamic range for transformer attention layers.
\end{itemize}

\subsection{RISC-V Matrix Extensions (RV-M / 2D Tiles)}
To natively support deep learning, the RISC-V ecosystem is standardizing the **RISC-V Matrix Extension (RV-M)**. Building upon the vector extension, RV-M organizes data into 2D matrix tiles (e.g., $m_0$ through $m_7$) and introduces dedicated matrix multiplication primitives that dispatch directly to hardware systolic arrays or 2D matrix execution engines.


\section{Comparative Synthesis of Domain-Specific Architectures}\label{sec:dsa_comparison}

Table \ref{tab:dsa_comparison} synthesizes the architectural trade-offs across all major paradigms of computing.

\begin{table}[htbp]
  \begin{center}
    \caption{Comparative Synthesis of General-Purpose and Domain-Specific Architectures.}
    \label{tab:dsa_comparison}
    \small
    \begin{tabular}{|l|l|l|l|l|l|}
      \hline
      \textbf{Dimension} & \textbf{General CPU} & \textbf{Vector Core} & \textbf{GPU (SIMT)} & \textbf{DSP} & \textbf{Tensor / TPU} \\ \hline \hline
      \textbf{Primary Target} & General serial & Scientific DLP & Massively parallel & Streaming signal & AI / GEMM \\ \hline
      \textbf{Programming}    & Sequential scalar & Vectorized VLA & Massively threaded & Fixed-point & Matrix / Graph \\ \hline
      \textbf{Hardware Concurrency} & ILP (1--8 issue) & DLP (lanes) & Thread / Warp & DSP MAC pipelines & 2D Systolic PEs \\ \hline
      \textbf{Control Overhead}     & Very High ($>80\%$) & Low ($<10\%$) & Low ($<15\%$) & Very Low ($<5\%$) & Negligible ($<2\%$) \\ \hline
      \textbf{Latency Strategy}     & Deep OOO / Cache & Pipelined lanes & Zero-cycle context switch & Deterministic pipe & Spatial pipelining \\ \hline
      \textbf{Memory Dataflow}      & Cache hierarchy & Strided / Gather & High-bandwidth HBM & Modified Harvard & Streaming / Systolic \\ \hline
      \textbf{Typical Precision}    & 32/64-bit int/FP & 32/64-bit FP/int & 16/32-bit FP & 16/32-bit fixed & INT8, FP8, BF16 \\ \hline
      \textbf{Energy Efficiency}    & Low ($\sim 1\times$) & High ($\sim 10\times$) & High ($\sim 15\times$) & Very High ($\sim 30\times$) & Peak ($\sim 50\text{--}100\times$) \\ \hline
    \end{tabular}
  \end{center}
\end{table}

By understanding these domain-specific architectural paradigms, the system designer can effectively partition computational workloads, combining the flexibility of general-purpose RISC-V cores with targeted hardware accelerators to overcome the power and memory walls of modern computing.
